\scanpage{349}
\mspage{13}

\uline{Proposition} : Soient $n,m\geq 2$, et considérons l’inclusion $U=\partial X\overset{i}{\hookrightarrow}X$ du bord du ``carré'' $X=\Phi_n\times\Phi_m$ dans $X$. C’est une immersion ouverte \inserted{$W$-}asphérique, de façon précise les $U_{/t}$ ($t\in\dot U=\overline U\setminus U$) sont isomorphes à $e$ (cat. ponctuelle) ou à
\diagram{\begin{tikzcd}[ampersand replacement=\&,row sep=tiny,column sep=small]
{} \& {} \& {}\arrow[dl]\\
\unsure{\Phi|\Psi}={} \& {}\\
{} \& {} \& {}\arrow[ul]
\end{tikzcd}}
(catégorie avec élément final), sauf \deleted{\unsure{pour}} les quatre éléments $t_i$ des quatre ``coins'' de $\dot U$, où c’est la catégorie
\diagram{\begin{tikzcd}[ampersand replacement=\&,row sep=tiny,column sep=small]
{} \& {} \& a \& a'\arrow[l]\\
\widetilde{\unsure{\Upsilon|Y}}={} \& o\arrow[ur]\arrow[dr] \& {}\\
{} \& {} \& b \& b'\arrow[l]
\end{tikzcd}}
\diagram{\begin{tikzcd}[ampersand replacement=\&,row sep=small,column sep=small]
o\arrow[r]\arrow[d] \& a\arrow[d] \& a'\arrow[l]\arrow[r]\arrow[d] \& {} \& \cdots\\
b\arrow[r] \& \bullet_{t_1} \& {}\arrow[l] \& {} \& \cdots\\
b'\arrow[u]\arrow[r]\arrow[d] \& {}\arrow[u] \& {}\arrow[u]\arrow[l] \& {} \& \cdots\\
{} \& \cdots \& \cdots \& {} \& \cdots
\end{tikzcd}}
n’ayant \deleted{\unsure{un} élément initial, mais} pas d’él\textsuperscript{t} final. Donc $U\xrightarrow{i}X$ n’est pas imm. ouverte de Dwyer, (mais $U\longrightarrow X\setminus\{t_1,t_2,t_3,t_4\}$ l’est...). Ceci posé, je dis que $i:U\longrightarrow X$ n’est pas une $W$-cofibration.

\uline{\unsure{Dém.}} Soit
\[
F\unsure{\subset|\lhook\joinrel\relbar}U
\]
l’inclusion dans $U=\partial X$ de la catégorie-\unsure{coin|coïn}
\diagram{\begin{tikzcd}[ampersand replacement=\&,row sep=small,column sep=small]
o\arrow[r]\arrow[d] \& a\\
b \& {}
\end{tikzcd}}
(cf figure), \inserted{isom. à $\Upsilon$.} \unsure{(}On voit aussitôt que c’est une imm. fermée de Dwyer (le plus petit ouvert de $U$ contenant $\Upsilon$ est précisément la catégorie $\widetilde\Upsilon=U_{/t_1}$ de tantôt), donc une $W$-cofibration. Si $U\unsure{\text{---}|\longrightarrow}X$ était une $W$-cofibration, il en serait de $=$ du composé $V\longrightarrow X$. On va voir que ce n’est pas le cas.

Soit
\diagram{\begin{tikzcd}[ampersand replacement=\&,row sep=small,column sep=small]
\overline V={} \& o\arrow[r]\arrow[d] \& a\arrow[d]\\
{} \& b\arrow[r] \& t_1
\end{tikzcd}}
l’adhérence de $V$ ds $X$,

% archive sheet number: 349
